\documentclass[11pt]{amsart}

\usepackage[a4paper,margin=2.6cm]{geometry}
\usepackage{amsmath,amssymb,amsthm,mathtools}
\usepackage{mathrsfs}
\usepackage{enumitem}
\usepackage{hyperref}

\newtheorem{theorem}{Theorem}
\newtheorem{proposition}[theorem]{Proposition}
\newtheorem{lemma}[theorem]{Lemma}
\newtheorem{corollary}[theorem]{Corollary}
\theoremstyle{remark}
\newtheorem{remark}[theorem]{Remark}

\newcommand{\C}{\mathbb C}
\newcommand{\R}{\mathbb R}
\newcommand{\D}{\mathcal D}
\newcommand{\supp}{\operatorname{supp}}
\newcommand{\1}{\mathbf 1}
\newcommand{\dd}{\,d}
\newcommand{\pb}{\partial_{\bar z}}

\title[Gaussian localization domains for Hermite eigenfunctions]{A detailed boundary-integral proof that Gaussian localization domains with Hermite eigenfunctions are disks or annuli}
\author{}
\date{}

\begin{document}

\maketitle

\begin{abstract}
We present a detailed write-up of a Bargmann-side argument showing that if a Hermite function $h_k$ is an eigenfunction of a localization operator with Gaussian window $h_0$ and symbol $\mathbf 1_U$, where $U\subset\C$ is bounded, connected, and of class $C^{1,\alpha}$, then $U$ is a disk or an annulus centered at the origin. The argument proceeds through the differentiated Bargmann identity, a compactly supported potential obtained by the Ehrenpreis--Malgrange divisibility step, one-phase regularity of the outer boundary, explicit boundary moment identities, a fully detailed steepest-descent lemma on the outer boundary, and a unique-continuation argument for the angular derivative. The only point treated as a black box is the Kinderlehrer--Nirenberg analyticity theorem for the one-phase free boundary.
\end{abstract}

\section{Statement of the result}

Let $L_U^{h_0}$ denote the localization operator with Gaussian window $h_0$ and symbol $\mathbf 1_U$.

\begin{theorem}\label{thm:main}
Let $U\subset\C$ be a bounded connected domain of class $C^{1,\alpha}$, and assume
\[
0\notin \partial U.
\]
Suppose that for some integer $k\ge 0$, the Hermite function $h_k$ is an eigenfunction of the localization operator $L_U^{h_0}$.

Assume the standard Bargmann reduction:
\begin{equation}\label{eq:exp-identity}
\int_U e^{\pi z w}\, |z|^{2k} e^{-\pi |z|^2}\, dA(z)=c_k
\qquad \forall w\in\C,
\end{equation}
for some constant $c_k\in\R$, and assume moreover the existence of a compactly supported distribution $u$ satisfying
\begin{equation}\label{eq:distribution-u}
\Delta u = f_k(|z|)\,\1_U - c_k\delta_0
\qquad\text{in }\D'(\C),
\end{equation}
where
\[
f_k(r):=r^{2k}e^{-\pi r^2}.
\]

Then $U$ is a disk or an annulus centered at the origin.
\end{theorem}

\begin{remark}
The differentiated identity \eqref{eq:exp-identity} is the standard consequence of differentiating the Bargmann eigenvalue identity $k$ times. The existence of a compactly supported $u$ satisfying \eqref{eq:distribution-u} is the output of the Ehrenpreis--Malgrange divisibility step applied to the compactly supported distribution
\[
f_k(|z|)\,\1_U-c_k\delta_0.
\]
The proof below starts from these two facts.
\end{remark}

\section{Preliminaries on the outer boundary}

Let $\Gamma_{\mathrm{out}}$ denote the outer boundary component of $U$, and let $E_\infty$ denote the unbounded connected component of $\C\setminus \overline U$.

\begin{lemma}\label{lem:outer-zero}
One has
\[
u\equiv 0
\qquad\text{in }E_\infty.
\]
\end{lemma}

\begin{proof}
Since $u$ is compactly supported, there exists $R_0>0$ such that
\[
\supp u \subset \overline{D_{R_0}(0)}.
\]
Hence
\[
u\equiv 0
\qquad\text{in }\C\setminus \overline{D_{R_0}(0)}.
\]
Because $U$ is bounded, after increasing $R_0$ if necessary we may assume
\[
\overline U \subset D_{R_0}(0).
\]
Now $E_\infty$ is connected, and in $E_\infty$ one has
\[
\Delta u=0
\]
in the distributional sense, because $E_\infty\cap U=\varnothing$ and $0\notin E_\infty$. Therefore $u$ is harmonic in $E_\infty$.

Since $E_\infty$ contains the exterior region $\C\setminus \overline{D_{R_0}(0)}$, where $u$ vanishes identically, unique continuation for harmonic functions yields
\[
u\equiv 0
\qquad\text{in }E_\infty.
\]
\end{proof}

\begin{lemma}\label{lem:local-one-phase}
Let $p\in \Gamma_{\mathrm{out}}$. Then there exists $r>0$ such that, with
\[
D:=D(p,r),
\]
one has:
\begin{enumerate}[label=\textnormal{(\roman*)}]
\item $0\notin \overline D$,
\item $D\cap \partial U = D\cap \Gamma_{\mathrm{out}}$,
\item
\[
\Delta u = f_k(|z|)
\qquad\text{in }D\cap U,
\]
\item
\[
u=0,\qquad \partial_\nu u=0
\qquad\text{on }D\cap \Gamma_{\mathrm{out}},
\]
where $\nu$ is the outer unit normal to $U$.
\end{enumerate}
\end{lemma}

\begin{proof}
Fix $p\in \Gamma_{\mathrm{out}}$. Since $0\notin \partial U$, there exists $r_1>0$ such that
\[
\overline{D(p,r_1)}\cap \{0\}=\varnothing.
\]
Since $\partial U$ is of class $C^{1,\alpha}$, every boundary point has a neighborhood in which the boundary is represented by a single $C^{1,\alpha}$ graph. Because $p$ lies on the outer boundary component, there exists $r_2>0$ such that
\[
D(p,r_2)\cap \partial U = D(p,r_2)\cap \Gamma_{\mathrm{out}}
\]
and
\[
D(p,r_2)\setminus \overline U \subset E_\infty.
\]
Set
\[
r:=\min\{r_1,r_2\},
\qquad
D:=D(p,r).
\]
Then (i) and (ii) hold by construction, and moreover
\[
D\setminus \overline U \subset E_\infty.
\]
By Lemma \ref{lem:outer-zero},
\[
u\equiv 0
\qquad\text{in }D\setminus \overline U.
\tag{$*$}
\]

Now let $\varphi\in C_c^\infty(D)$. Since $0\notin \overline D$, the singular term in \eqref{eq:distribution-u} does not contribute, and so
\[
\langle \Delta u,\varphi\rangle
=
\int_U f_k(|z|)\varphi(z)\,dA(z)
=
\int_{D\cap U} f_k(|z|)\varphi(z)\,dA(z).
\]
Thus
\[
\Delta u = f_k(|z|)
\qquad\text{in }\D'(D\cap U),
\]
which proves (iii). Since the right-hand side is smooth in $\overline D$, elliptic regularity implies $u\in C^\infty(D\cap U)$.

Because $u$ vanishes identically on the exterior side by ($*$), continuity up to the boundary gives
\[
u=0
\qquad\text{on }D\cap \Gamma_{\mathrm{out}}.
\]

It remains to show that the interior normal derivative also vanishes there. Let again $\varphi\in C_c^\infty(D)$. Green's formula in the smooth domain $D\cap U$ gives
\[
\int_{D\cap U} u\,\Delta\varphi\,dA
=
\int_{D\cap U} (\Delta u)\varphi\,dA
-
\int_{D\cap \Gamma_{\mathrm{out}}} \partial_\nu u\, \varphi\, ds,
\]
because $u=0$ on $D\cap \Gamma_{\mathrm{out}}$. On the other hand, the distributional identity gives
\[
\int_{D\cap U} u\,\Delta\varphi\,dA
=
\langle \Delta u,\varphi\rangle
=
\int_{D\cap U} f_k(|z|)\varphi\,dA.
\]
Since $\Delta u=f_k(|z|)$ in $D\cap U$, comparison yields
\[
\int_{D\cap \Gamma_{\mathrm{out}}} \partial_\nu u\, \varphi\, ds = 0
\qquad\forall\,\varphi\in C_c^\infty(D).
\]
Hence
\[
\partial_\nu u=0
\qquad\text{on }D\cap \Gamma_{\mathrm{out}}.
\]
This proves (iv).
\end{proof}

\begin{proposition}\label{prop:outer-analytic}
The outer boundary component $\Gamma_{\mathrm{out}}$ is a real-analytic Jordan curve.
\end{proposition}

\begin{proof}
Fix $p\in \Gamma_{\mathrm{out}}$, and let $D=D(p,r)$ be as in Lemma \ref{lem:local-one-phase}. In $D\cap U$, the equation reads
\[
\Delta u = f_k(|z|),
\]
and on $D\cap \Gamma_{\mathrm{out}}$ one has
\[
u=0,\qquad \partial_\nu u=0.
\]
Since the tangential derivative of the trace is automatically zero (the trace is identically zero), we also have
\[
\nabla u=0
\qquad\text{on }D\cap \Gamma_{\mathrm{out}}.
\]

Now $f_k(|z|)$ is real-analytic in a neighborhood of $\overline D$ because $0\notin \overline D$, and
\[
f_k(r)=r^{2k}e^{-\pi r^2}>0
\qquad (r>0).
\]
Thus near $D\cap \Gamma_{\mathrm{out}}$ we are in the setting of a one-phase free-boundary problem with analytic, strictly positive forcing:
\[
\Delta u=f_k(|z|)\quad\text{in }D\cap U,
\qquad
u=0,\ \nabla u=0\quad\text{on }D\cap \Gamma_{\mathrm{out}}.
\]
By the classical theorem of Kinderlehrer and Nirenberg, a $C^{1,\alpha}$ free boundary for such a problem is real-analytic. Therefore $D\cap \Gamma_{\mathrm{out}}$ is a real-analytic arc.

Since $p$ was arbitrary, $\Gamma_{\mathrm{out}}$ is globally a real-analytic Jordan curve.
\end{proof}

\section{Boundary integral identities}

Define, for $s\ge 0$,
\begin{equation}\label{eq:Hk-def}
H_k(s):=s^{-(k+1)}\int_0^s t^k e^{-\pi t}\,dt.
\end{equation}
By l'H\^opital's rule, $H_k$ extends smoothly across $s=0$, and it is real-analytic on $(0,\infty)$.

\begin{lemma}\label{lem:Hk-ode}
For every $s\ge 0$ one has
\begin{equation}\label{eq:Hk-identity}
(k+1)H_k(s)+sH_k'(s)=e^{-\pi s}.
\end{equation}
Consequently, for every integer $n\ge 0$,
\begin{equation}\label{eq:pb-primitive}
\pb\!\Big(z^{n+k}\bar z^{\,k+1}H_k(|z|^2)\Big)
=
z^n |z|^{2k}e^{-\pi|z|^2}.
\end{equation}
\end{lemma}

\begin{proof}
Differentiate \eqref{eq:Hk-def}:
\[
H_k'(s)
=
-(k+1)s^{-(k+2)}\int_0^s t^k e^{-\pi t}\,dt
+s^{-(k+1)}s^k e^{-\pi s}.
\]
Multiplying by $s$ and adding $(k+1)H_k(s)$ gives \eqref{eq:Hk-identity}.

Now set
\[
F(z,\bar z):=z^{n+k}\bar z^{\,k+1}H_k(|z|^2).
\]
Using $\pb(|z|^2)=z$, we compute
\begin{align*}
\pb F
&=
z^{n+k}\Big((k+1)\bar z^k H_k(|z|^2)+\bar z^{k+1}H_k'(|z|^2)\,z\Big) \\
&=
z^{n+k}\bar z^k\Big((k+1)H_k(|z|^2)+|z|^2H_k'(|z|^2)\Big) \\
&=
z^{n+k}\bar z^k e^{-\pi|z|^2}
=
z^n |z|^{2k}e^{-\pi|z|^2}.
\end{align*}
This proves \eqref{eq:pb-primitive}.
\end{proof}

\begin{lemma}\label{lem:boundary-moments}
For every integer $n\ge 1$,
\begin{equation}\label{eq:bulk-moment}
\int_U z^n |z|^{2k}e^{-\pi|z|^2}\,dA(z)=0
\end{equation}
and
\begin{equation}\label{eq:boundary-moment}
\int_{\partial U} z^{n+k}\bar z^{\,k+1}H_k(|z|^2)\,dz = 0.
\end{equation}
\end{lemma}

\begin{proof}
Expand the exponential in \eqref{eq:exp-identity}:
\[
e^{\pi zw}=\sum_{n=0}^\infty \frac{\pi^n}{n!} z^n w^n.
\]
Since the right-hand side of \eqref{eq:exp-identity} is constant in $w$, comparison of coefficients gives \eqref{eq:bulk-moment} for all $n\ge 1$.

Now apply the Cauchy--Green formula
\[
\int_U \pb F\, dA=\frac{1}{2i}\int_{\partial U} F(z,\bar z)\,dz
\]
to
\[
F(z,\bar z)=z^{n+k}\bar z^{\,k+1}H_k(|z|^2).
\]
Using \eqref{eq:pb-primitive} and \eqref{eq:bulk-moment}, we obtain
\[
0
=
\int_U z^n |z|^{2k}e^{-\pi|z|^2}\,dA(z)
=
\frac{1}{2i}\int_{\partial U} z^{n+k}\bar z^{\,k+1}H_k(|z|^2)\,dz,
\]
which is \eqref{eq:boundary-moment}.
\end{proof}

\section{Detailed asymptotic analysis of the outer boundary}

In this section we prove, in full detail, that the outer boundary component
$\Gamma_{\mathrm{out}}$ is a circle centered at the origin.

\subsection{The componentwise oscillatory integrals}

Let
\[
\partial U=\Gamma_{\mathrm{out}}\cup\Gamma_1\cup\cdots\cup\Gamma_m
\]
be the decomposition of the boundary into connected components, each endowed
with the orientation induced by $U$.

For the outer boundary we fix a real-analytic parametrization
\[
\gamma:[0,2\pi]\to \Gamma_{\mathrm{out}}
\]
such that
\[
\gamma'(t)\neq 0
\qquad\text{for all }t\in[0,2\pi].
\]
Likewise, for each inner component $\Gamma_j$ we fix a $C^1$ parametrization
\[
\gamma_j:[0,2\pi]\to \Gamma_j.
\]

By Lemma~\ref{lem:boundary-moments}, for every integer $n\ge 1$,
\[
\int_{\partial U} z^{n+k}\bar z^{\,k+1}H_k(|z|^2)\,dz=0.
\]
Splitting this over the boundary components yields
\begin{equation}\label{eq:split-boundary}
I_n^{\mathrm{out}}+\sum_{j=1}^m I_n^{(j)}=0,
\qquad n\ge 1,
\end{equation}
where
\begin{equation}\label{eq:I-out}
I_n^{\mathrm{out}}
:=
\int_0^{2\pi}
\gamma(t)^{n+k}\overline{\gamma(t)}^{\,k+1}
H_k(|\gamma(t)|^2)\gamma'(t)\,dt,
\end{equation}
and
\begin{equation}\label{eq:I-inner}
I_n^{(j)}
:=
\int_0^{2\pi}
\gamma_j(t)^{n+k}\overline{\gamma_j(t)}^{\,k+1}
H_k(|\gamma_j(t)|^2)\gamma_j'(t)\,dt.
\end{equation}

For the outer boundary it is convenient to write
\begin{equation}\label{eq:A-def}
A(t):=\gamma(t)^k\overline{\gamma(t)}^{\,k+1}
H_k(|\gamma(t)|^2)\gamma'(t),
\end{equation}
so that
\begin{equation}\label{eq:I-out-short}
I_n^{\mathrm{out}}=\int_0^{2\pi}\gamma(t)^n A(t)\,dt.
\end{equation}

\begin{lemma}\label{lem:A-nonzero}
The function $A$ is real-analytic and never vanishes on $[0,2\pi]$.
\end{lemma}

\begin{proof}
Since $\Gamma_{\mathrm{out}}$ is a real-analytic Jordan curve and $0\notin\Gamma_{\mathrm{out}}$, the functions
\[
\gamma,\qquad \overline{\gamma},\qquad |\gamma|^2,\qquad H_k(|\gamma|^2),\qquad \gamma'
\]
are all real-analytic on $[0,2\pi]$; therefore $A$ is real-analytic.

Now $\gamma(t)\neq 0$ for every $t$, because $0\notin \Gamma_{\mathrm{out}}$; moreover $\gamma'(t)\neq 0$ by assumption. Finally,
\[
H_k(s)=s^{-(k+1)}\int_0^s t^k e^{-\pi t}\,dt>0
\qquad (s>0),
\]
since the integrand is positive. Therefore every factor in \eqref{eq:A-def} is nonzero, and hence
\[
A(t)\neq 0
\qquad\text{for all }t\in[0,2\pi].
\]
\end{proof}

\subsection{A detailed asymptotic lemma for one analytic component}

\begin{lemma}[Detailed componentwise asymptotic expansion]\label{lem:component-asymptotic-detailed}
Let $\gamma:[0,2\pi]\to\C$ be a real-analytic $2\pi$-periodic map such that
\[
\gamma'(t)\neq 0
\qquad\text{for all }t\in[0,2\pi],
\]
and let $A:[0,2\pi]\to\C$ be a real-analytic function. For $n\ge 1$, define
\[
I_n:=\int_0^{2\pi}\gamma(t)^nA(t)\,dt.
\]

Set
\[
R(t):=|\gamma(t)|,
\qquad
R_*:=\max_{t\in[0,2\pi]}R(t).
\]
Assume that $R$ is not constant. Then the set of maximizers
\[
\mathcal M:=\{t\in[0,2\pi]:R(t)=R_*\}
\]
is finite. Write
\[
\mathcal M=\{t_1,\dots,t_N\},
\qquad
\gamma(t_j)=R_*e^{i\theta_j}.
\]

Assume moreover that
\[
A(t_j)\neq 0
\qquad (j=1,\dots,N).
\]

For each $j$ there exist:
\begin{itemize}
\item an integer $m_j\ge 1$,
\item a real number $a_j>0$,
\end{itemize}
such that
\begin{equation}\label{eq:R-expansion}
R(t)=R_* - a_j (t-t_j)^{2m_j} + O\!\bigl((t-t_j)^{2m_j+1}\bigr)
\qquad\text{as }t\to t_j.
\end{equation}

Let
\[
m_*:=\max_{1\le j\le N}m_j,
\qquad
J_*:=\{j:\,m_j=m_*\}.
\]
Then there exist nonzero constants $c_j\in\C$ for $j\in J_*$ such that
\begin{equation}\label{eq:full-asymptotic}
I_n
=
R_*^n n^{-1/(2m_*)}
\left(
\sum_{j\in J_*} c_j e^{in\theta_j}
\right)
+
o\!\left(R_*^n n^{-1/(2m_*)}\right)
\qquad (n\to\infty).
\end{equation}
\end{lemma}

\begin{proof}
We divide the proof into several steps.

\medskip

\noindent\textbf{Step 1: finiteness of the maximum set.}
Since $R$ is real-analytic and not constant, the level set
\[
\{t\in[0,2\pi]:R(t)=R_*\}
\]
cannot have an accumulation point unless $R\equiv R_*$ identically. Since we assume $R$ nonconstant, $\mathcal M$ is finite.

\medskip

\noindent\textbf{Step 2: local expansions near a maximum.}
Fix $j\in\{1,\dots,N\}$. Since $R$ is real-analytic and has a maximum at $t_j$, all odd derivatives of $R$ up to the first nonzero derivative vanish, and the first nonzero derivative is of even order and negative. Hence there exist an integer $m_j\ge 1$ and a positive number $a_j$ such that \eqref{eq:R-expansion} holds.

Since $\gamma(t_j)\neq 0$, there exists a real-analytic branch of the argument near $t_j$; that is, we may write
\[
\gamma(t)=R(t)e^{i\varphi(t)}
\]
for $t$ near $t_j$, where $\varphi$ is real-analytic and
\[
\varphi(t_j)=\theta_j.
\]
Expanding at $t_j$ gives
\[
\varphi(t)=\theta_j+b_j(t-t_j)+O((t-t_j)^2)
\]
for some real number $b_j$.

Likewise, since $A$ is real-analytic,
\[
A(t)=A(t_j)+O(t-t_j)
\qquad\text{as }t\to t_j.
\]

\medskip

\noindent\textbf{Step 3: decomposition of the integral.}
Choose $\delta>0$ so small that the closed intervals
\[
J_j:=[t_j-\delta,t_j+\delta]
\]
are pairwise disjoint and contain no other points of $\mathcal M$.

Decompose
\[
I_n=\sum_{j=1}^N I_n^{(j)}+E_n,
\]
where
\[
I_n^{(j)}:=\int_{J_j}\gamma(t)^nA(t)\,dt,
\qquad
E_n:=\int_{[0,2\pi]\setminus \bigcup_{j=1}^N J_j}\gamma(t)^nA(t)\,dt.
\]

Since $R(t)<R_*$ on the compact set
\[
[0,2\pi]\setminus \bigcup_{j=1}^N J_j,
\]
there exists $\delta_0>0$ such that
\[
R(t)\le R_*-\delta_0
\qquad\text{for all }t\in [0,2\pi]\setminus \bigcup_{j=1}^N J_j.
\]
Hence
\[
|E_n|
\le
2\pi \|A\|_{L^\infty}(R_*-\delta_0)^n.
\]
Since
\[
(R_*-\delta_0)^n=o(R_*^n n^{-M})
\]
for every fixed $M>0$, we conclude that
\begin{equation}\label{eq:E-small}
E_n=o(R_*^n n^{-M})
\qquad\text{for every }M>0.
\end{equation}

\medskip

\noindent\textbf{Step 4: local analysis near one maximum point.}
Fix $j$. On $J_j$, write
\[
\gamma(t)^n=R(t)^n e^{in\varphi(t)}.
\]
Since
\[
R(t)=R_*\left(1-\frac{a_j}{R_*}(t-t_j)^{2m_j}+O((t-t_j)^{2m_j+1})\right),
\]
taking logarithms gives
\[
\log\!\frac{R(t)}{R_*}
=
-\frac{a_j}{R_*}(t-t_j)^{2m_j}
+O((t-t_j)^{2m_j+1}).
\]
Thus
\[
\left(\frac{R(t)}{R_*}\right)^n
=
\exp\!\left(
-\frac{n a_j}{R_*}(t-t_j)^{2m_j}
+
O\!\bigl(n|t-t_j|^{2m_j+1}\bigr)
\right).
\]
Combining this with the expansion of $\varphi$ and of $A$, we obtain
\begin{align*}
\gamma(t)^nA(t)
&=
R_*^n e^{in\theta_j}
\exp\!\left(
-\frac{n a_j}{R_*}(t-t_j)^{2m_j}
\right)
\bigl(A(t_j)+o(1)\bigr),
\end{align*}
uniformly on $J_j$ as $n\to\infty$.

Therefore
\[
I_n^{(j)}
=
R_*^n e^{in\theta_j}
\int_{J_j}
\exp\!\left(
-\frac{n a_j}{R_*}(t-t_j)^{2m_j}
\right)
\bigl(A(t_j)+o(1)\bigr)\,dt.
\]

Now set
\[
t-t_j=n^{-1/(2m_j)}y.
\]
Then
\[
dt=n^{-1/(2m_j)}dy,
\]
and
\[
\exp\!\left(
-\frac{n a_j}{R_*}(t-t_j)^{2m_j}
\right)
=
\exp\!\left(-\frac{a_j}{R_*}y^{2m_j}\right).
\]
Thus
\[
I_n^{(j)}
=
R_*^n e^{in\theta_j} n^{-1/(2m_j)}
\int_{-\delta n^{1/(2m_j)}}^{\delta n^{1/(2m_j)}}
\exp\!\left(-\frac{a_j}{R_*}y^{2m_j}\right)
\bigl(A(t_j)+o(1)\bigr)\,dy.
\]
By dominated convergence,
\[
\int_{-\delta n^{1/(2m_j)}}^{\delta n^{1/(2m_j)}}
\exp\!\left(-\frac{a_j}{R_*}y^{2m_j}\right)
\bigl(A(t_j)+o(1)\bigr)\,dy
\to
A(t_j)\int_{\R}\exp\!\left(-\frac{a_j}{R_*}y^{2m_j}\right)\,dy.
\]
Since $A(t_j)\neq 0$ and the integral is strictly positive, the limit is a nonzero complex number. Hence
\begin{equation}\label{eq:Ij-asymptotic}
I_n^{(j)}
=
R_*^n e^{in\theta_j}
\left(
c_j n^{-1/(2m_j)}
+
o\!\left(n^{-1/(2m_j)}\right)
\right),
\end{equation}
where
\[
c_j
:=
A(t_j)\int_{\R}\exp\!\left(-\frac{a_j}{R_*}y^{2m_j}\right)\,dy
\neq 0.
\]

\medskip

\noindent\textbf{Step 5: assemble the contributions.}
Summing \eqref{eq:Ij-asymptotic} over $j=1,\dots,N$ and using \eqref{eq:E-small}, we obtain
\[
I_n
=
R_*^n
\sum_{j=1}^N e^{in\theta_j}
\left(
c_j n^{-1/(2m_j)}
+
o\!\left(n^{-1/(2m_j)}\right)
\right)
+
o(R_*^n n^{-M})
\]
for every fixed $M>0$.

Among the exponents $n^{-1/(2m_j)}$, the slowest decay occurs precisely for those $j$ with $m_j=m_*$. Therefore the indices in $J_*$ provide the leading term, and we arrive at
\[
I_n
=
R_*^n n^{-1/(2m_*)}
\left(
\sum_{j\in J_*} c_j e^{in\theta_j}
\right)
+
o\!\left(R_*^n n^{-1/(2m_*)}\right).
\]
This is \eqref{eq:full-asymptotic}.
\end{proof}

\begin{lemma}[A finite exponential sum cannot tend to zero]\label{lem:exp-sum-detailed}
Let $\lambda_1,\dots,\lambda_L\in\C$ satisfy
\[
|\lambda_\ell|=1,
\qquad
\lambda_\ell\neq \lambda_m \text{ for }\ell\neq m.
\]
Let $c_1,\dots,c_L\in\C$. If
\[
\sum_{\ell=1}^L c_\ell \lambda_\ell^n \longrightarrow 0
\qquad (n\to\infty),
\]
then
\[
c_1=\cdots=c_L=0.
\]
\end{lemma}

\begin{proof}
Fix $j\in\{1,\dots,L\}$. Multiply the relation by $\overline{\lambda_j}^{\,n}$:
\[
\overline{\lambda_j}^{\,n}\sum_{\ell=1}^L c_\ell\lambda_\ell^n
=
c_j+\sum_{\ell\neq j} c_\ell (\lambda_\ell\overline{\lambda_j})^n.
\]
Average from $n=1$ to $N$:
\[
\frac1N\sum_{n=1}^N \overline{\lambda_j}^{\,n}\sum_{\ell=1}^L c_\ell\lambda_\ell^n
=
c_j
+
\sum_{\ell\neq j} c_\ell
\frac1N\sum_{n=1}^N (\lambda_\ell\overline{\lambda_j})^n.
\]
Since $\lambda_\ell\neq \lambda_j$, we have
\[
\lambda_\ell\overline{\lambda_j}\neq 1.
\]
Hence
\[
\frac1N\sum_{n=1}^N (\lambda_\ell\overline{\lambda_j})^n \to 0
\qquad (N\to\infty)
\]
for every $\ell\neq j$.

On the other hand, the hypothesis implies
\[
\overline{\lambda_j}^{\,n}\sum_{\ell=1}^L c_\ell\lambda_\ell^n \to 0,
\]
because $|\lambda_j|=1$. Therefore the Ces\`aro means on the left-hand side also tend to $0$. Passing to the limit gives
\[
c_j=0.
\]
Since $j$ was arbitrary, all coefficients vanish.
\end{proof}

\begin{proposition}\label{prop:outer-circle-detailed}
The outer boundary component $\Gamma_{\mathrm{out}}$ is a circle centered at the origin.
\end{proposition}

\begin{proof}
Let
\[
\gamma:[0,2\pi]\to \Gamma_{\mathrm{out}}
\]
be a real-analytic parametrization of the outer boundary. Define
\[
I_n^{\mathrm{out}}
:=
\int_0^{2\pi}
\gamma(t)^{n+k}\overline{\gamma(t)}^{\,k+1}
H_k(|\gamma(t)|^2)\gamma'(t)\,dt.
\tag{4.4}
\]
By Lemma \ref{lem:A-nonzero}, this is of the form
\[
I_n^{\mathrm{out}}=\int_0^{2\pi}\gamma(t)^nA(t)\,dt
\]
with $A$ real-analytic and nonvanishing.

Set
\[
R_{\mathrm{out}}(t):=|\gamma(t)|,
\qquad
R_*:=\max_{t\in[0,2\pi]} R_{\mathrm{out}}(t).
\]

Assume for contradiction that $\Gamma_{\mathrm{out}}$ is not a circle centered at the origin. Then $R_{\mathrm{out}}$ is nonconstant. Hence Lemma \ref{lem:component-asymptotic-detailed} applies, yielding
\begin{equation}\label{eq:outer-asymptotic}
I_n^{\mathrm{out}}
=
R_*^n n^{-1/(2m_*)}
\left(
\sum_{\ell=1}^L c_\ell e^{in\theta_\ell}
\right)
+
o\!\left(R_*^n n^{-1/(2m_*)}\right),
\end{equation}
for some $m_*\ge 1$, nonzero coefficients $c_\ell$, and distinct phases $e^{i\theta_\ell}$.

We now compare this with the contributions of the inner boundary components. For each $j=1,\dots,m$, define
\[
R_j^*:=\max_{z\in \Gamma_j}|z|.
\]
Since $\Gamma_j$ lies strictly inside the bounded Jordan domain enclosed by $\Gamma_{\mathrm{out}}$, the subharmonic function $|z|^2$ attains its maximum on that closed region on the outer boundary, and the maximum is strict on $\Gamma_j$. Therefore
\[
R_j^*<R_*
\qquad (j=1,\dots,m).
\tag{4.5}
\]

Now fix $j\in\{1,\dots,m\}$. Since $\Gamma_j$ is compact and the integrand in \eqref{eq:I-inner} is continuous, there exists $C_j>0$ such that
\[
|I_n^{(j)}|
\le
C_j (R_j^*)^n
\qquad (n\ge 1).
\]
Because $R_j^*<R_*$, this implies
\[
I_n^{(j)}=o\!\left(R_*^n n^{-1/(2m_*)}\right).
\]
Summing over all inner components, we obtain
\[
\sum_{j=1}^m I_n^{(j)}
=
o\!\left(R_*^n n^{-1/(2m_*)}\right).
\tag{4.6}
\]

But by the boundary identity \eqref{eq:split-boundary},
\[
I_n^{\mathrm{out}}+\sum_{j=1}^m I_n^{(j)}=0.
\]
Combining \eqref{eq:outer-asymptotic} and \eqref{eq:4.6}, and dividing by $R_*^n n^{-1/(2m_*)}$, gives
\[
\sum_{\ell=1}^L c_\ell e^{in\theta_\ell}\to 0.
\]
By Lemma \ref{lem:exp-sum-detailed}, this is impossible unless every $c_\ell$ is zero, contradicting the asymptotic lemma.

Therefore our assumption was false. Hence $R_{\mathrm{out}}$ is constant, which means that $\Gamma_{\mathrm{out}}$ is a circle centered at the origin.
\end{proof}

\section{Propagation of radiality}

\begin{lemma}\label{lem:remove-point}
If $U\subset\C$ is open and connected, then $U\setminus\{0\}$ is connected.
\end{lemma}

\begin{proof}
Open connected subsets of $\R^2$ are path connected. Let $x,y\in U\setminus\{0\}$. Choose a continuous path in $U$ joining $x$ to $y$. If the path does not meet $0$, we are done. Otherwise, since $U$ is open, a small disk around $0$ is contained in $U$, and the portion of the path passing through that disk can be replaced by a circular arc avoiding $0$. Thus $x$ and $y$ can still be joined inside $U\setminus\{0\}$.
\end{proof}

\begin{proposition}\label{prop:u-radial}
The potential $u$ is radial in $U\setminus\{0\}$.
\end{proposition}

\begin{proof}
By Proposition \ref{prop:outer-circle-detailed}, the outer boundary is the circle
\[
\Gamma_{\mathrm{out}}=\{|z|=R\}
\]
for some $R>0$.

Define the angular derivative
\[
L:=x\partial_y-y\partial_x=\partial_\theta.
\]
Since $L$ commutes with $\Delta$ and annihilates every radial distribution, applying $L$ to \eqref{eq:distribution-u} gives
\[
\Delta (Lu)=0
\qquad\text{in }U\setminus\{0\}.
\]
Set
\[
w:=Lu.
\]
Then $w$ is harmonic in $U\setminus\{0\}$.

Because the outer boundary is a circle, compactness of the inner boundary components implies that there exists $\varepsilon>0$ such that the collar
\[
\mathcal A:=\{R-\varepsilon<|z|<R\}
\]
satisfies
\[
\mathcal A\subset U.
\]
Indeed, the union of the inner boundary components is compact and disjoint from $\Gamma_{\mathrm{out}}$.

On the outer boundary circle $|z|=R$, Lemma \ref{lem:local-one-phase} gives
\[
u=0,\qquad \partial_r u=0.
\]
Therefore
\[
w=\partial_\theta u=0
\qquad\text{on }|z|=R,
\]
and also
\[
\partial_r w=\partial_r\partial_\theta u=\partial_\theta\partial_r u=0
\qquad\text{on }|z|=R.
\]

Now $w$ is harmonic in the annulus $\mathcal A$, so it admits the Fourier expansion
\[
w(r,\theta)
=
a_0+b_0\log r
+\sum_{n\ge 1}(a_n r^n+b_n r^{-n})\cos(n\theta)
+\sum_{n\ge 1}(c_n r^n+d_n r^{-n})\sin(n\theta).
\]
The boundary conditions
\[
w(R,\theta)=0,\qquad \partial_r w(R,\theta)=0
\]
force every coefficient in this expansion to vanish. Hence
\[
w\equiv 0
\qquad\text{in }\mathcal A.
\]

By Lemma \ref{lem:remove-point}, the set $U\setminus\{0\}$ is connected. Since $w$ is harmonic there and vanishes in the nonempty open subset $\mathcal A$, unique continuation gives
\[
w\equiv 0
\qquad\text{in }U\setminus\{0\}.
\]
Thus
\[
\partial_\theta u=0
\qquad\text{in }U\setminus\{0\},
\]
so $u$ is radial in $U\setminus\{0\}$.
\end{proof}

\section{Proof of the theorem}

\begin{proof}[Proof of Theorem \ref{thm:main}]
By Proposition \ref{prop:u-radial}, the distribution $u$ is radial in $U\setminus\{0\}$. Therefore $\Delta u$ is a radial distribution.

From \eqref{eq:distribution-u},
\[
\Delta u = f_k(|z|)\,\1_U-c_k\delta_0.
\]
The point mass $\delta_0$ is radial, hence the distribution
\[
f_k(|z|)\,\1_U
\]
must also be radial.

Because
\[
f_k(r)=r^{2k}e^{-\pi r^2}>0
\qquad (r>0),
\]
it follows that for each fixed $r>0$ the slice
\[
\{\theta\in\R/2\pi\Z:\ re^{i\theta}\in U\}
\]
has characteristic function equal almost everywhere to a constant. Since this slice is open in the circle, it must be either empty or the whole circle. Thus $U$ is rotationally invariant about the origin.

A bounded connected rotationally invariant planar domain is either a disk or an annulus. Therefore $U$ is a disk or an annulus centered at the origin.
\end{proof}

\begin{thebibliography}{99}

\bibitem{KN}
D. Kinderlehrer and L. Nirenberg,
\emph{Regularity in free boundary problems},
Ann. Scuola Norm. Sup. Pisa Cl. Sci. (4) \textbf{4} (1977), 373--391.

\bibitem{Ramos}
J. P. G. Ramos,
\emph{Eigenfunctions of time-frequency localisation operators and free-boundary problems in potential theory},
preprint.

\end{thebibliography}

\end{document}